\chapter{Control Flow}
\label{ch:control-flow}

Until now our programs have run straight through, top to bottom, doing the same
thing every time. Real programs make \emph{decisions}: do this when the user is
old enough, do that otherwise. The tool for choosing between paths is the
\code{if} statement, and it turns the boolean expressions of Chapter~3 into
action.

\section{The \texttt{if} statement}

An \code{if} runs a block of code only when a condition is true. The condition is
any expression of type \code{bool}; the block opens with a colon and closes with
\code{end}:

\begin{salam}
func main:
    age : int = 20
    if age >= 18:
        println "You may enter."
    end
    println "Done."
end
\end{salam}

\begin{progout}
You may enter.
Done.
\end{progout}

If \code{age} were \code{15}, the condition \code{age >= 18} would be
\code{false}, the indented block would be skipped entirely, and only
\code{Done.} would print. The structure is always the same: \code{if}, a
condition, a colon, the body, and \code{end}.

\section{\texttt{else}: the other path}

Often you want one thing to happen when the condition holds and a \emph{different}
thing otherwise. That is what \code{else} is for:

\begin{salam}
func main:
    age : int = 15
    if age >= 18:
        println "You may enter."
    else:
        println "Sorry, too young."
    end
end
\end{salam}

\begin{progout}
Sorry, too young.
\end{progout}

Exactly one of the two blocks runs never both, never neither. Think of it as a
fork in the road: the condition decides which way to go.

\section{\texttt{else if}: choosing among many}

When there are more than two possibilities, chain them with \code{else if}. Salam
checks each condition in order and runs the block belonging to the \emph{first}
one that is true; if none match, the final \code{else} (if present) runs.

\begin{tip}
Notice that Salam does not repeat the word \code{if} after \code{else}. Put a
condition directly after \code{else} and it behaves as \code{else if}; leave
\code{else} with no condition and it is the final, catch-all branch.
\end{tip}

\begin{salam}
func main:
    score : int = 83
    if score >= 90:
        println "Grade: A"
    else score >= 80:
        println "Grade: B"
    else score >= 70:
        println "Grade: C"
    else:
        println "Grade: F"
    end
end
\end{salam}

\begin{progout}
Grade: B
\end{progout}

The order matters. \code{83} satisfies \code{>= 80}, so we print \code{B} and
stop the remaining branches are not even tested. Because the checks run top
to bottom, arrange them from the most specific or highest threshold downward, as
here.

\section{The inline form}

When a branch's body is a single short statement, you may write it on the same
line as the colon. This keeps tidy little decisions compact:

\begin{salam}
func main:
    n : int = 7
    if n % 2 == 0: println "even"
    else:          println "odd"
    end
end
\end{salam}

\begin{progout}
odd
\end{progout}

The \code{end} is still required; only the line breaks change. Many of the
example programs that ship with Salam use this form to line up short branches
neatly, like a small table. Use it when it improves readability, and the
multi-line form when a body grows beyond a single statement.

\section{Conditions are just boolean expressions}

There is nothing special about the condition in an \code{if} it is any
expression that produces a \code{bool}, exactly the kind you built in Chapter~3.
That means you can combine tests with \code{\&\&}, \code{||}, and \code{!}:

\begin{salam}
func main:
    age : int = 20
    has_ticket : bool = true
    if age >= 18 && has_ticket:
        println "Welcome to the show."
    else:
        println "Entry denied."
    end
end
\end{salam}

\begin{progout}
Welcome to the show.
\end{progout}

You can also store a condition in a \code{bool} variable first and test that ---
sometimes a well-named boolean makes the intent far clearer than a long condition
crammed into the \code{if}:

\begin{salam}
func main:
    age : int = 20
    has_ticket : bool = true
    let_in : auto = age >= 18 && has_ticket
    if let_in:
        println "Welcome to the show."
    end
end
\end{salam}

\begin{warn}
Salam's conditions must be genuine \code{bool} values. Unlike some languages, you
cannot use a number as a shortcut for ``truth'' there is no rule that
\code{0} means false and everything else means true. Write the comparison out:
\code{count != 0} rather than just \code{count}. This keeps conditions explicit
and removes a whole category of subtle bugs.
\end{warn}

\section{Nesting}

The body of an \code{if} is an ordinary block, so it may contain another
\code{if}. This is called \emph{nesting}, and it lets you ask a follow-up
question only after a first one has been answered:

\begin{salam}
func main:
    logged_in : bool = true
    is_admin : bool = false
    if logged_in:
        if is_admin:
            println "Welcome, administrator."
        else:
            println "Welcome, user."
        end
    else:
        println "Please log in."
    end
end
\end{salam}

\begin{progout}
Welcome, user.
\end{progout}

\begin{tip}
Deep nesting quickly becomes hard to follow. Two habits keep it shallow: combine
related tests with \code{\&\&} instead of nesting two \code{if}s, and once you
reach Chapter~6 use an early \code{ret} to handle a special case and leave the
main path un-indented. If you find yourself four levels deep, that is usually a
sign to restructure.
\end{tip}

\section{Decisions inside functions}

Conditions become especially powerful inside functions (the subject of
Chapter~6), where a branch can \emph{return} a result. Here is a function that
classifies a score, using \code{else if} with the inline \code{ret} form:

\begin{salam}
func grade(score : int) : str:
    if score >= 90: ret "A"
    else score >= 80: ret "B"
    else score >= 70: ret "C"
    else: ret "F"
    end
end

func main:
    println grade(95), grade(83), grade(71), grade(50)
end
\end{salam}

\begin{progout}
A B C F
\end{progout}

Each branch hands back a different string, and the first matching one wins. We
will return to functions properly soon; for now, notice how naturally
\code{if}/\code{else if} expresses ``pick the right answer.''

\section{In Persian}

Conditionals read just as naturally in Persian. The keywords are \code{\fa{اگر}}
for \code{if}, \code{\fa{وگرنه}} for \code{else}, and \code{\fa{تمام}} for
\code{end}:

\begin{salamfa}
تابع اصلی:
    سن : صحیح = 15
    اگر سن >= 18:
        چاپ "بزرگسال"
    وگرنه:
        چاپ "کم‌سن"
    تمام
تمام
\end{salamfa}

An \code{else if} chain is written the same way as in English: put a condition
right after \code{\fa{وگرنه}}, with no separate word for \code{if}.

\section{Chapter summary}

\begin{itemize}[leftmargin=1.4em]
  \item \code{if cond:} runs a block only when \code{cond} (a \code{bool}) is
        true; close it with \code{end}.
  \item \code{else} provides the alternative path; exactly one branch runs.
  \item \code{else if} chains several conditions; the first true one wins.
  \item A single-statement branch may be written inline after the colon.
  \item Conditions are ordinary boolean expressions and must be genuine
        \code{bool}s numbers are not truthy.
  \item \code{if}s may be nested, but prefer \code{\&\&} and early returns to keep
        nesting shallow.
\end{itemize}

\section{Exercises}

\exercise Write a program that reads a fixed integer variable \code{temp} and
prints \code{"freezing"} if it is below 0, \code{"cold"} if below 15,
\code{"warm"} if below 30, and \code{"hot"} otherwise.

\exercise Using a single \code{if} with \code{\&\&}, print \code{"valid"} only
when a variable \code{hour} is between 0 and 23 inclusive.

\exercise Write a function \code{sign(n int) str} that returns \code{"negative"},
\code{"zero"}, or \code{"positive"}, and test it with three values.

\exercise Rewrite the grade example so the inline branches become multi-line
blocks that each also print a short message. Does the output change?

\exercise Translate the freezing/cold/warm/hot program into Persian.
